\documentclass[UTF8,zihao=-4]{ctexart}

\usepackage[a4paper,margin=2.3cm]{geometry}
\usepackage{array}
\usepackage{booktabs}
\usepackage{longtable}
\usepackage{hyperref}
\usepackage{xcolor}
\usepackage{enumitem}
\usepackage{fancyhdr}
\usepackage{lastpage}

\hypersetup{
  colorlinks=true,
  linkcolor=black,
  urlcolor=blue
}

\setlength{\parindent}{2em}
\setlength{\parskip}{0.35em}
\renewcommand{\arraystretch}{1.25}
\setlist[itemize]{leftmargin=2.2em,itemsep=0.2em,topsep=0.2em}
\pagestyle{fancy}
\fancyhf{}
\fancyfoot[C]{第 \thepage 页 / 共 \pageref*{LastPage} 页}
\renewcommand{\headrulewidth}{0pt}
\renewcommand{\footrulewidth}{0pt}

\title{\heiti 托福阅读 Class 1 课堂笔记}
\author{}
\date{}

\begin{document}
\maketitle
\thispagestyle{fancy}

本节课主要学习托福阅读中常见的词根、词缀和易混词。通过词根词缀可以帮助我们推测单词含义，提高阅读中理解生词的能力。

\section*{一、词根词缀与重点单词}

\texttt{-phot-} 通常与“光”有关。例如 \texttt{photosynthesizing} 表示“能进行光合作用的”，\texttt{euphotic} 表示“光照充足的、透光的”。在自然科学类阅读中，这一词根常与光照、光合作用和水域透光层等内容相关。

\texttt{eu-} 是前缀，通常表示“好的、优越的”。例如：\texttt{eugenics} 表示“优生学”，\texttt{euphonious} 表示“悦耳的”，\texttt{euphemism} 表示“委婉语”，\texttt{euthanasia} 表示“安乐死”，\texttt{euphoria} 表示“兴奋感、欣快感”，\texttt{eulogy} 表示“悼词、颂词”。

\texttt{-phyt-} 通常与“植物”有关。例如：\texttt{phytoplankton} 表示“浮游植物”，\texttt{lithophyte} 是“石生植物”，\texttt{hydrophyte} 是“水生植物”，\texttt{xerophyte} 是“旱生植物”，\texttt{epiphyte} 是“附生植物”，\texttt{cryophyte} 是“寒生植物”。需要注意的是，\texttt{zooplankton} 表示“浮游动物”，其中 \texttt{zoo-} 与动物有关。

\texttt{supercalifragilisticexpialidocious} 可理解为“极好的、奇妙无比的”。课堂解释为“用精致的美丽来弥补可教育性的缺失”。其中，\texttt{super-} 表示“超级的”；\texttt{cali-} 与“美”有关，如 \texttt{calisthenics} 表示“健美操”，\texttt{calligraphy} 表示“书法”；\texttt{fragil-} 表示“脆弱的、精细的”，如 \texttt{fragment} 表示“碎片”，\texttt{fraction} 表示“部分、分数”，\texttt{frail} 表示“虚弱的、脆弱的”；\texttt{expia-} 表示“补偿、弥补”，如 \texttt{expiate} 表示“赎罪、弥补”，\texttt{expiation} 表示“赎罪、补偿”；\texttt{doc-} 表示“教”，如 \texttt{doctor} 表示“博士、医生”，\texttt{document} 表示“文件”，\texttt{docile} 表示“温顺的、易受教的”，\texttt{doctrine} 表示“教义、学说”，\texttt{docent} 表示“讲师、讲解员”。

\texttt{sanct-} 表示“神圣的”。相关单词包括：\texttt{sanctuary}，意为“避难所、庇护所、圣所”；\texttt{sanctum}，意为“圣所、私室”；\texttt{sanctimonious}，意为“假装虔诚的、道貌岸然的、伪善的”。注意：手写笔记中出现的 \texttt{sanctorium} 容易与 \texttt{sanatorium} 混淆，常见正式词 \texttt{sanatorium} 表示“疗养院”，与“治疗、恢复健康”有关，不属于 \texttt{sanct-} 这一组。

\texttt{cur / curs} 表示“跑、流动”。相关单词包括：\texttt{cursor}，意为“光标、游标”；\texttt{cursory}，意为“草率的、粗略的”；\texttt{cursive script}，意为“草书、连笔字”；\texttt{discursive article}，意为“东拉西扯的文章、散漫的文章”；\texttt{excursion}，意为“短途旅行”；\texttt{incursion}，意为“入侵”；\texttt{precursor}，意为“先驱、前兆、前体物质”；\texttt{recursive function}，意为“递归函数”。此外，\texttt{expedition} 表示“远征、探险、考察”，\texttt{forerunner} 表示“先驱、前身、预兆”，\texttt{predecessor} 表示“前任、前身”。

\texttt{path- / pathy / pathos} 可以与“感情、感受”有关，也可以与“疾病”有关。\texttt{apathy} 表示“冷漠、漠不关心”，其中 \texttt{a-} 表示“没有”；\texttt{empathy} 表示“共情、同理心”，\texttt{create empathy} 表示“引起共鸣”，\texttt{empathetic} 表示“有同理心的”；\texttt{sympathy} 表示“同情”，其中 \texttt{sym- / syn-} 表示“一起、共同”；\texttt{telepathy} 表示“心灵感应”，其中 \texttt{tele-} 表示“远的”；\texttt{pathetic} 表示“可怜的、令人同情的”；\texttt{pathology} 表示“病理学”；\texttt{pathogen} 表示“病原体”；\texttt{cardiopathy} 表示“心脏病”。

\texttt{mal-} 表示“不好的、坏的”。例如：\texttt{malevolence} 表示“恶意”，\texttt{malfunction} 表示“故障、失灵”，\texttt{malnutrition} 表示“营养不良”，\texttt{malignant tumor} 表示“恶性肿瘤”。与之相对，\texttt{benign} 表示“良性的”。

\texttt{trans-} 表示“穿过、转变”。例如：\texttt{transit} 表示“中转、运输”，\texttt{transition} 表示“过渡、转变”，\texttt{transient} 和 \texttt{transitory} 都可表示“短暂的、暂时的”。

\texttt{-escent} 表示“逐渐……的”。例如：\texttt{adolescent} 表示“青少年”，\texttt{evanescent} 表示“逐渐消失的、短暂的”，\texttt{quiescent} 表示“静止的、休眠的”。相关词 \texttt{descent} 表示“下降”。

\section*{二、易混词与补充表达}

\begin{longtable}{>{\raggedright\arraybackslash}p{4.5cm} >{\raggedright\arraybackslash}p{10cm}}
\toprule
\textbf{单词或表达} & \textbf{含义} \\
\midrule
\endfirsthead
\toprule
\textbf{单词或表达} & \textbf{含义} \\
\midrule
\endhead
\texttt{appreciable} & 可观的，值得注意的 \\
\texttt{buoyancy} & 浮力；也可表示乐观、轻快 \\
\texttt{sustenance} & 食物，营养；维持生命的东西 \\
\texttt{catastrophic} & 灾难性的 \\
\texttt{hysterical} & 歇斯底里的，情绪失控的 \\
\texttt{vertical} & 垂直的 \\
\texttt{prerequisite} & 先决条件 \\
\texttt{furnish} & 提供，供应 \\
\texttt{frugal} & 节约的，节俭的 \\
\texttt{parsimonious} & 过分节俭的，吝啬的 \\
\texttt{sanguine} & 脸色红润的；乐观的，满怀希望的 \\
\texttt{trivial} & 琐碎的，不重要的 \\
\texttt{futile} & 徒劳的，无用的 \\
\texttt{eminent} & 杰出的，著名的 \\
\texttt{preeminent} & 卓越的，出类拔萃的 \\
\texttt{imminent} & 即将发生的，迫近的 \\
\texttt{intermittent} & 间歇的，断断续续的 \\
\texttt{pertinent} & 相关的，切题的 \\
\texttt{prominent} & 突出的，显著的，著名的 \\
\texttt{renounce} & 放弃，宣布拒绝；否认与某事有关 \\
\texttt{pronounced} & 显著的，明显的 \\
\texttt{enunciate} & 清楚地发音；清楚表达 \\
\texttt{denounce} & 谴责，公开指责 \\
\texttt{perfunctory} & 敷衍的，草率的 \\
\texttt{sentimental} & 多愁善感的，感伤的 \\
\texttt{sensation} & 感觉；轰动，引起轰动的人或事 \\
\texttt{sensory} & 感官的，感觉的 \\
\texttt{sensible} & 明智的，合理的 \\
\texttt{adhere to} & 坚持，遵守；粘附于 \\
\texttt{adherent} & 支持者，拥护者 \\
\texttt{adhesive} & 黏合的；黏合剂 \\
\texttt{inherent} & 固有的，内在的 \\
\texttt{coherent} & 连贯的，条理清楚的 \\
\texttt{coherence} & 连贯性，一致性 \\
\texttt{cohesive} & 有凝聚力的；黏着的 \\
\bottomrule
\end{longtable}

修辞学中常见的三个概念是 \texttt{pathos}、\texttt{logos} 和 \texttt{ethos}。\texttt{pathos} 指“感染力、情感说服”，\texttt{logos} 指“理性、逻辑说服”，\texttt{ethos} 指“可信度、品格说服”。

经济产业分类中，\texttt{the primary sector} 表示“第一产业”，\texttt{the secondary sector} 表示“第二产业”，\texttt{the tertiary sector} 表示“第三产业”。

\end{document}
